\chapter{Prédateur/proie}
\mysetref{fox-rabbit}

Existe t-il une méthode de transformation systématique\footnote{%
  Pouvant à terme devenir un algorithme ou au  moins une heuristique.} %
d'une description différentielle continue d'un système en équation
cellulaire? %
Confronté au problème de la découverte d'une règle de transition
codant un comportement\footnote{%
  Le problème de la réduction d'un comportement à une dynamique, ou
  plus gé\-né\-ra\-le\-ment la détermination des observables de
  l'espace-temps, nous amène à une utilisation plus métaphorique
  de termes désignant les données initiales ou résultantes
  d'une expérience cellulaire.}, %
la démarche usuelle repose plus sur la recherche intuitive d'une
description phénoménologique locale, que sur l'utilisation
combinatoire d'un répertoire de primitives doté d'un comportement
prévisible\footnote{%
  Voir~\cite{Richards90} pour une recherche de règles à partir
  d'espaces de configurations, utilisant un algorithme génétique pour
  la recherche dans l'espace des règles, mais concernant des AC
  probabilistes.}. %
Partant d'un comportement global donné, cette démarche intuitive
postule l'existence d'un comportement local associé à un observable
global. %
Une règle locale est obtenue par la simplification symbolique
d'une description analytique. %
Comme exemple de l'utilisation de cette méthode empirique, prenons
un problème classique de dynamique des populations. %
Nous avons choisi un système déjà fort complexe, permettant une
formulation cellulaire:~un système proie/prédateur à deux espèces. %
Le squelette de la fonction de transition s'obtient aisément par la
description littérale des comportements individuels supposés d'une
population naturelle:

\begin{itemize}
\item Les individus naissent, durent et meurent.
\item Les proies meurent en présence des prédateurs.
\end{itemize}

Un bit d'existence et $n$ bits de longévité sont alloués dans l'espace
cellulaires à chaque espèce\footnote{% 
  Nous avons utilisé des compteurs sur 3 bits.}. %
La naissance d'un individu dépend de la densité de la population locale. %
Les paramètres facteurs de croissance et de décroissance de chaque
population ne sont donc pas fixés de manière absolue, mais résultent de
la conjonction du compteur de longévité et des seuils locaux
déclenchant une naissance. %
A l'observation, les variations de population résultantes présentent
l'aspect cyclique corrélé attendu, mais le répertoire des
comportements accessibles via les paramètres de la fonction de
transition ou par des modifications du squelette de
l'automate\footnote{%
  Par exemple en imposant la présence de proies dans le voisinage
  d'une cellule, en plus de la présence de prédateurs, comme une
  condition nécessaire à la naissance d'un prédateur.} %
est difficilement comparable à l'expression différentielle classique
de ce problème.

\mytwoplotVfull{rabbit-gnans-1}{rabbit-gnans-2}%
{\asoft{gnans} et Lokta-Volterra}%
{Les équations~\ref{rabbit-1} et~\ref{rabbit-2} calculées par~\asoft{gnans}}%
{Deux solutions périodiques obtenues avec \protect\mysoftcite{gnans} pour
  l'équation du modèle prédateur-proie de Lokta-Volterra.}

La figure~\mygetplot{rabbit-gnans-1+rabbit-gnans-2}
est obtenue avec \mysoftcite{gnans} pour l'équation
du modèle prédateur-proie de Lokta-Volterra. %
D'après~\cite{Murray90}:

\begin{eqnarray}
  \frac{dN}{dt} & = N(a-bP), \label{rabbit-1}\\
  \frac{dP}{dt} & = P(cN-d), \label{rabbit-2}
\end{eqnarray}

Avec a, b, c et d constantes positives. %
Les équations~\ref{rabbit-1} et~\ref{rabbit-2} sont spécifiées tel
que à~\asoft{gnans}:

\footnotesize
\begin{verbatim}
        CONTINUOUS TIME SYSTEM lokta;
        
        AT TIME t:
        d(n) = n * (a - b * p);
        d(p) = p * (c * n - dd);
\end{verbatim}
\normalsize

Les valeurs des variables et des paramètres sont, pour la
solution~1\footnote{%
  Et les conditions opérationnelles:~%
  timestep = 1, %
  stoptime = 48, %
  algo = RKF45.}:

\begin{footnotesize}
\begin{verbatim}
        STATE n = 2;
        STATE p = 1.99;
        TIME t;
        
        PARAMETER a = 4;
        PARAMETER b = 3;
        PARAMETER c = 2;
        PARAMETER dd = 5;
\end{verbatim}
\end{footnotesize}

Et pour la solution~2\footnote{%
  Timestep = 1, stoptime = 200, algo = Ode.}:

\footnotesize
\begin{verbatim}
        STATE n = 1;
        STATE p = 1;
        TIME t;
        
        PARAMETER a = .8;
        PARAMETER b = .1;
        PARAMETER c = .1;
        PARAMETER dd = .1;
\end{verbatim}
\normalsize

\mytwoplotVfull{rabbit-1}{rabbit-1-512}%
{\arule{Fox-Rabbit} Version 1}%
{Une tentative transposition cellulaire}%
{Évolution des populations proie-prédateurs obtenue avec la règle\\
  \arule{Fox-Rabbit} version~1}

La figure~\mygetplot{rabbit-1+rabbit-1-512} montre l'évolution des
populations proie-prédateurs obtenue avec la règle \arule{Fox-Rabbit}
version~1, pour 100~générations sur un espace taille \asquare{128}~(haut)
et \asquare{512}~(bas). %
Dans les deux cas, l'espace initial est de densité aléatoire\,\undemi~%
sur le plan \avar{fox} et sur le plan \avar{rabbit}.

\medskip
\large
\arule{Fox-Rabbit} version~1:
\normalsize

\footnotesize
\begin{verbatim}
        fox_sum = W0 + N0 + S0 + E0;
        rabbit_sum = W1 + N1 + S1 + E1;
        ofox = ifox;
        orabbit = irabbit;

        if (ifox && irabbit) {
                orabbit = 0;
        }
        if (ifox) {
                ofox_life = ifox_life - 1;
                ofox = ofox_life > 0;
        } 
        if (irabbit && !ifox) {
                orabbit_life = irabbit_life - 1;
                orabbit = orabbit_life > 0;
        }
        if (!ifox && fox_sum >= 2 && rabbit_sum >= 1) {
                ofox = 1;
                ofox_life = 3;
        }
        if (!irabbit && rabbit_sum >= 2) {
                orabbit = 1;
                orabbit_life = 3;
        }
\end{verbatim}
\normalsize

\mytwoplotVfull{rabbit-2}{rabbit-3}%
{}%
{\arule{Fox-Rabbit} Versions~2 et~3}%
{Nous considérons la variété des comportements dynamique obtenue pour
  les perturbations apportés à la règle initiale\footnote{%
    Compte tenu de la taille du voisinage~(2~plans von Neumann de
    1~bit chacun et 6~bits locaux).} comme un indicateur de l'intérêt
  d'une recherche empirique d'une typologie métaphorique des
  dynamiques cellulaires, \ie d'une tentative de classification des
  scénarios utilisés pour établir des expériences cellulaires.}

La figure~\mygetplot{rabbit-2+rabbit-3} montre l'évolution des
populations proie-prédateurs obtenue avec la règle \arule{Fox-Rabbit}
version 2~(haut) et 3~(bas). %
Pour~200 et 400~générations sur un espace taille~\asquare{128}, même
espace initial que pour la figure~\mygetplot{rabbit-1+rabbit-1-512}.

\medskip
\large
\arule{Fox-Rabbit} version~2\footnote{%
  La dernière condition de
  création d'un rabbit est suspecte!}:
\normalsize

\footnotesize
\begin{verbatim}
        if (ifox && irabbit) {
                ofox_life = ifox_life - 1;
                ofox = ofox_life > 0;
                orabbit = 0;
        }
        if (ifox && !irabbit) {
                ofox_life = ifox_life - 3;
                ofox = ofox_life > 0;
        } 
        if (irabbit && !ifox) {
                orabbit_life = irabbit_life - 1;
                orabbit = orabbit_life > 0;
        }
        if (!ifox && fox_sum == 2 && rabbit_sum >= 1) {
                ofox = 1;
                ofox_life = 7;
        }
        if (!irabbit && (rabbit_sum >= 2 || rabbit_sum <= 3)) {
                orabbit = 1;
                orabbit_life = 2;
        }
\end{verbatim}
\normalsize

\large
\arule{Fox-Rabbit} version~3:
\normalsize

\footnotesize
\begin{verbatim}
        if (ifox && irabbit) {
                if (fox_sum > rabbit_sum) {
                        ofox_life = ifox_life - 1;
                        ofox = ofox_life > 0;
                        orabbit = 0;
                } else if (fox_sum < rabbit_sum) {
                        orabbit_life = irabbit_life - 1;
                        orabbit = orabbit_life > 0;
                        ofox = 0;
                }
        }
        if (ifox && !irabbit) {
                ofox_life = ifox_life - 2;
                ofox = ofox_life > 0;
        }
        if (irabbit && !ifox) {
                orabbit_life = irabbit_life - 2;
                orabbit = orabbit_life > 0;
        }
        if (!ifox && fox_sum == 2) {
                ofox = 1;
                ofox_life = 7;
        }
        if (!irabbit && rabbit_sum >= 2 && fox_sum >= 1) {
                orabbit = 1;
                orabbit_life = 3;
        }
\end{verbatim}
\normalsize

%%
%% Local Variables:
%% mode: latex
%% mode: auto-fill
%% iso-accents-minor-mode: nil
%% TeX-master: "top"
%% TeX-command-default: "mytex"
%% Local IspellParsing: "latex-mode ~latin1"
%% Local IspellPersDict: "~/.ispell_lisp"
%% LocalWords: "fox-rabbit t-il g\'e ra ph\'enom\'enologique rabbit-gnans gnans AT"
%% LocalWords: "Lokta-Volterra pr\'edateur-proie CONTINUOUS TIME SYSTEM lokta dd"
%% LocalWords: "timestep stoptime algo RKF STATE PARAMETER rabbit Fox-Rabbit sum"
%% LocalWords: "proie-pr\'edateurs ofox ifox orabbit irabbit von Neumann bug\'e else"
%% End:
